\documentclass[10pt,a4paper]{amsart}
\usepackage[margin=19mm]{geometry}
\usepackage{amsmath,amssymb,mathtools}
\usepackage[scr=rsfs,cal=euler]{mathalfa}
\usepackage{xfrac}
\usepackage[unicode,hidelinks,bookmarks=false]{hyperref}
\newcommand{\ie}{i.e.\ }
\newcommand{\cf}{cf.\ }
\newcommand{\resp}{respectively\ }
\newcommand{\Subsection}{\subsection}
\providecommand{\nobreakdash}{\nobreak\hbox{-}}
\setlength{\emergencystretch}{2em}
\multlinegap0pt
\usepackage{amssymb}
\usepackage{xcolor}
\newtheorem{thm}{Theorem}
\DeclareMathOperator{\Ad}{Ad}
\newcommand{\SL}{\textrm{SL}}
\newcommand{\Tr}{\mathrm{Tr}}
\newcommand{\bC}{\mathbb{C}}
\newcommand{\cZ}{\mathcal{Z}}
\newcommand{\fg}{\mathfrak{g}}
\newcommand{\pr}{\mathrm{pr}}
\title{Geometry of the tt*-Toda equations \\ I: universal centralizer \\ and symplectic groupoids}
\author{Martin A. Guest, Nan-Kuo Ho}
\date{Source: arXiv:2604.04358v1 (CC BY 4.0)}
\begin{document}
\maketitle

\noindent\emph{Source and licence.} Geometry of the tt*-Toda equations \\ I: universal centralizer \\ and symplectic groupoids. Version arXiv:2604.04358v1 (CC BY 4.0), distributed by arXiv under the Creative Commons Attribution 4.0 International licence, http://creativecommons.org/licenses/by/4.0/, as recorded in the arXiv licence metadata for this exact version and shown on the abstract page https://arxiv.org/abs/2604.04358v1. Full licence notice: \texttt{licenses/arxiv-2604.04358-CC-BY-4.0.txt}.\par\smallskip

\noindent\textit{Editorial scope.} Counted unit: the article's thm thm:main (source0.tex lines 2092--2094). The article's complete original proof of it is delivered with it (source0.tex lines 2096--2206, one \texttt{proof} environment of 111 lines). No mathematics is written, repaired or re-derived: the statement and the proof are the article's own lines, in the article's own notation, and no retained line is substituted or re-flowed.  Retained uncounted as the source-local context the proof needs: lines 1866--1868; lines 1934--1936.  Nothing was omitted from any retained span, so the package carries no omission record.  No retained line contains a citation, so the wrapper carries no bibliography and no bibliography entry.  No retained line contains a figure environment, an image-inclusion command or drawing code, so no image is delivered and the package declares no asset dependency.  Editorial formatting only, in addition: an AMS \texttt{amsart} preamble that restates the article's own declarations and macros used by the retained lines, namely the environments \texttt{thm} on separate counters, together with the standard packages those lines need.  The retained environments print in this document as Theorem 1 in order of appearance, whereas the article prints the same items with the numbering of its own full document.  The complete native source of the article as distributed in the arXiv e-print for this exact version is bundled unchanged as \texttt{sources/197879/source0.tex}.
\begin{equation}\label{eq:quasi2form}
\omega_{(g,a)}=\tfrac{1}{2}[(\Ad_a\pr_1^*\theta^L,\pr_1^*\theta^L)+(\pr_1^*\theta^L,\pr^*_2(\theta^L+\theta^R))]
\end{equation}

\begin{thm}\label{thm:universalsymplectic}
Let $G$ be a complex, semisimple, and simply connected Lie group. Then the universal centralizer of a Steinberg cross section $\cZ\rightrightarrows \Sigma$ is a holomorphic symplectic groupoid with the holomorphic symplectic form $i^*\omega$, where $i:\cZ\hookrightarrow G\times G$ is the natural inclusion and $\omega$ is defined by equation (\ref{eq:quasi2form}).
\end{thm}

\begin{thm}\label{thm:main}
(i) $\sigma^*\omega_\bC=\omega_\bC$ along the units $\epsilon(\Sigma)$. (ii) $\theta^*\omega_{\bC}=-\overline{\omega}_\bC$ along the units $\epsilon(\Sigma)$.
\end{thm}

\begin{proof}
We shall use a similar idea to that in the proof of Theorem \ref{thm:universalsymplectic}. Recall the natural splitting of the tangent space $T_{(id_G,M)}\cZ=T_{(id_G,M)}\epsilon(\Sigma)\oplus T_{(id_G,M)}(s^{-1}(M))$ along the units, and write $u=u^H+u^F$ for the associated decomposition of $u\in T_{(id_G,M)}\cZ$.

We begin by expressing the tangent vectors as velocity vectors of suitable curves.
A tangent vector $u\in T_{(id_G,M)}\cZ$ can be expressed by
$$u=\tfrac{d}{dt}|_{t=0}(r_E(t),L_M(r_M(t)))\in T_{(id_G,M)}\cZ,$$
where $r_E(t)$ and $r_M(t)$ are curves in $G$ with $r_E(0)=id_G=r_M(0)$, such that the curve $L_M(r_M(t))\in\Sigma$ and $r_E(t)\in G$, which satisfy the commuting property for some neighborhood of $0$. Here $L_M$ denotes the left group multiplication by $M$. Hence the tangent vector
$$u=(r'_E(0),(L_M)_*(r'_M(0)) )=( (r'_E(0))^L(id_G), (r'_M(0))^L(M) ),$$
has the decomposition $u^H=\frac{d}{dt}|_{t=0}(id_G,L_M(r_M(t)))=(0,(r'_M(0))^L(M))$ and $u^F=u-u^H=( (r'_E(0))^L(id_G),0 )$. We shall express $u^F=\frac{d}{dt}|_{t=0}(\tilde{r}_E(t),M)$ for some $\tilde{r}_E(t)\in G_M$, where $\tilde{r}_E(0)=id_G$, $\tilde{r}'_E(0)=r'_E(0) \in \fg_M$.
Here the notation $\xi^L(g)$ means the left invariant vector field of $\xi\in\fg$ evaluated at $g\in G$.

Similarly, we express another tangent vector $v$ by
$$v=\tfrac{d}{dt}|_{t=0}(\delta_E(t),L_M(\delta_M(t)))\in T_{(id_G,M)}\cZ,$$ with the same condition.
Then $v^H=\frac{d}{dt}|_{t=0}(id_G,L_M(\delta_M(t)))=(0,(\delta'_M(0))^L(M))$ and $v^F=v-v^H=( (\delta'_E(0))^L(id_G),0 )$. We shall express $v^F=\frac{d}{dt}|_{t=0}(\tilde{\delta}_E(t),M)$ for some $\tilde{\delta}_E(t)\in G_M$, where $\tilde{\delta}_E(0)=id_G$, $\tilde{\delta}'_E(0)=\delta'_E(0) \in \fg_M$.

To simplify notation, the holomorphic symplectic form $\omega_\bC$ will be denoted by $\omega$ for the rest of the proof. Recall from the calculation in the proof of Theorem \ref{thm:universalsymplectic} that
$$\omega_{(id_G,M)}(u,v)=\omega_{(id_G,M)}(u^H,v^F)+\omega_{(id_G,M)}(u^F,v^H)$$
because $\omega_{(id_G,M)}(u^H,v^H)=0$ and $\omega_{(id_G,M)}(u^F,v^F)=0$. Thus we shall only need to compute the cross terms.

{\bf (i) The map $\sigma$}

Since $(\sigma,\sigma_0)$ is a groupoid morphism, direct calculation shows that $(\sigma_*u)^H=\sigma_*(u^H)$ and $(\sigma_*u)^F=\sigma_*(u^F)$. 
Thus $\omega_{\sigma(id_G,M)}(\sigma_*u,\sigma_*v)=\omega_{\sigma(id_G,M)}(\sigma_*(u^H),\sigma_*(v^F))
+\omega_{\sigma(id_G,M)}(\sigma_*(u^F),\sigma_*(v^H)),$  and  we only need to find the formula for $\sigma_*(u^H)$ and $\sigma_*(u^F)$.

Denote the tangent vectors by $\xi_1=\tilde{r}_E'(0),~\eta_1=\tilde{\delta}_E'(0)\in \fg_M$, and $\xi_2=r_M'(0), ~\eta_2=\delta_M'(0)\in\fg$, so $u=( \xi_1^L(id_G),\xi_2^L(M) )$ and $v=(\eta_1^L(id_G),\eta^L_2(M))$.
Then using the curves $\tilde{r}_E(t)$ and $\tilde{\delta}_E(t)$, we obtain
\begin{eqnarray*}
&&\sigma_*(u^F)=\tfrac{d}{dt}|_{t=0} \left( F_\sigma(M)(\tilde{r}_E(t)^{-T})F_\sigma(M)^{-1},  \sigma_o(M) \right)=\left( (-\Ad_{F_\sigma(M)} \xi_1^T)^L (id_G), 0 \right) \\
&&\sigma_*(v^F)=\tfrac{d}{dt}|_{t=0} \left( F_\sigma(M)(\tilde{\delta}_E(t)^{-T})F_\sigma(M)^{-1},  \sigma_o(M) \right)=\left( (-\Ad_{F_\sigma(M)}\eta_1^T)^L(id_G),0 \right).
\end{eqnarray*}
Next, $\sigma_*(u^H)=\frac{d}{dt}|_{t=0} (id_G, F_\sigma(L_M(r_M(t))) (L_M(r_M(t)))^{-T} F_\sigma(L_M(r_M(t)))^{-1} \in T_{(id_G,\sigma_o(M))}\epsilon(\Sigma)$, so
\begin{eqnarray*}
&&\sigma_*(u^H)=\left( 0, \left(\Ad_{F_\sigma(M) M^T F_\sigma(M)^{-1}}( \tfrac{d}{dt}|_{t=0}F_\sigma(L_M(r_M(t)))F_\sigma(M)^{-1}) \right)^L (\sigma_o(M)) \right) \\&&
+ \left( 0, \left( -\Ad_{F_\sigma(M)}\xi_2^T \right)^L(\sigma_o(M) ) \right) + \left( 0,\left( -\tfrac{d}{dt}|_{t=0}F_\sigma(L_M(r_M(t)))F_\sigma(M)^{-1} \right)^L( \sigma_o(M) ) \right),
\end{eqnarray*}
and
\begin{eqnarray*}
&&\sigma_*(v^H)=\left( 0, \left(\Ad_{F_\sigma(M)M^TF_\sigma(M)^{-1}}( \tfrac{d}{dt}|_{t=0}F_\sigma(L_M(\delta_M(t)))F_\sigma(M)^{-1}) \right)^L (\sigma_o(M) ) \right) \\&&
+ \left( 0, \left(-\Ad_{F_\sigma(M)}\eta_2^T\right)^L(\sigma_o(M) )\right) + \left( 0, \left(-\tfrac{d}{dt}|_{t=0}F_\sigma(L_M(\delta_M(t)))F_\sigma(M)^{-1} \right)^L( \sigma_o(M) ) \right).
\end{eqnarray*}

Calculating the value of the symplectic form with these curves, we obtain
\begin{eqnarray*}
\lefteqn{2 \omega_{(id_G,\sigma_o(M))}(\sigma_*(u^H),\sigma_*(v^F)) }\\
&&=\left( \eta_1^T,-\xi_2^T-\Ad_{M^{-T}}(\xi_2^T) \right)+\left( \Ad_{M^{-T}} (\eta_1^T) -\Ad_{M^T} (\eta_1^T), \Ad_{F_\sigma(M)^{-1}}( \tfrac{d}{dt}|_{t=0} F_\sigma(L_Mr_M(t))F_\sigma(M)^{-1}) \right)
\end{eqnarray*}
Since $\eta_1\in\fg_M$, $\Ad_M\eta_1=\eta_1=\Ad_{M^{-1}}\eta_1$ and $\Ad_{M^{-T}}(\eta_1^T)=\eta_1^T=\Ad_{M^T}\eta^T_1$. 
Hence the second term in the last equation vanishes, and we have obtained
\begin{equation}\label{eq:sigma_uHvF}
2\omega_{(id_G,\sigma_o(M) )}(\sigma_*(u^H),\sigma_*(v^F))=(\eta_1^T,-\xi_2^T-\Ad_{M^{-T}}(\xi_2^T) )=2(\eta_1^T, -\xi_2^T)
\end{equation}

A similar calculation gives 
\begin{equation}\label{eq:sigma_uFvH}
2\omega_{(id_G, \sigma_o(M) )}(\sigma_*(u^F),\sigma_*(v^H))=-(\xi_1^T,-\eta_2^T-\Ad_{M^{-T}}(\eta_2^T) )=-2(\xi_1^T,-\eta_2^T)
\end{equation}

On the other hand, we know from the calculation in the proof of Theorem \ref{thm:universalsymplectic} that
\begin{equation}\label{eq:uHvF}
2\omega_{(id_G,M)}(u^H,v^F)=(-\eta_1,\xi_2+\Ad_M\xi_2)=2(-\eta_1,\ \xi_2)
\end{equation}
and
\begin{equation}\label{eq:uFvH}
2\omega_{(id_G,M)}(u^F,v^H)=(\xi_1,\eta_2+\Ad_M\eta_2)=2(\xi_1,\ \eta_2)
\end{equation}
Since $(X,Y)=\Tr(XY)$ for $\SL(n,\bC)$, we have
$$(\eta_1^T,-\xi_2^T )=-\Tr(\xi_2\eta_1)=-(\eta_1,\xi_2), \quad\mbox{and}\quad (-\xi_1^T,-\eta_2^T)=\Tr(\xi_1\eta_2)=(\xi_1,\eta_2),$$
so equation (\ref{eq:sigma_uHvF}) is equation (\ref{eq:uHvF}), and equation (\ref{eq:sigma_uFvH}) is equation (\ref{eq:uFvH}). 
Since $u$, $v$ are arbitrary, we conclude that $\sigma^*\omega=\omega$ along $\epsilon(\Sigma)$ in $\cZ$.

{\bf (ii) The map $\theta$}

For the same reason as before, we only need to calculate $\theta_*(u^H)$ and $\theta_*(u^F)$.
We have
\[(\theta^*\omega)_{(id_G,M)}(u,v)
=\omega_{\theta(id_G,M)}(\theta_*(u^H),\theta_*(v^F))+\omega_{\theta(id_G,M)}(\theta_*(u^F),\theta_*(v^H)).
\]

Using the same description of the tangent vector $u$ and $v$ and arguing as in (i), we have
\[\theta_*(u^F)=( (\Ad_{F_\theta(M)} \overline{\xi}_1)^L (id_G), 0 )\quad \mbox{and} \quad
\theta_*(v^F)=( (\Ad_{F_\theta(M)}\overline{\eta}_1)^L(id_G),0 ).\]
for the vertical part.

For the horizontal part we have:
\begin{eqnarray*}
&&\theta_*(u^H)=\left( 0, \left( \Ad_{F_\theta(M) \overline{M} F_\theta(M)^{-1}}( \tfrac{d}{dt}|_{t=0}F_\theta(L_M(r_M(t)))F_\theta(M)^{-1}) \right)^L (\theta_o(M) ) \right) \\&&
+ \left( 0, \left( -\Ad_{F_\theta(M)}\overline{\xi}_2 \right)^L(\theta_o(M) ) \right) + \left( 0,\left( -\tfrac{d}{dt}|_{t=0}F_\theta(L_M(r_M(t)))F_\theta(M)^{-1} \right)^L( \theta_o(M) ) \right),
\end{eqnarray*}
and
\begin{eqnarray*}
&&\theta_*(v^H)=\left( 0, \left( \Ad_{F_\theta(M) \overline{M} F_\theta(M)^{-1}}( \tfrac{d}{dt}|_{t=0}F_\theta(L_M(\delta_M(t)))F_\theta(M)^{-1}) \right)^L (\theta_o(M) ) \right) \\&&
+ \left( 0, \left( -\Ad_{F_\theta(M)}\overline{\eta}_2 \right)^L(\theta_o(M) ) \right) + \left( 0,\left( -\tfrac{d}{dt}|_{t=0}F_\theta(L_M(\delta_M(t)))F_\theta(M)^{-1} \right)^L( \theta_o(M) ) \right).
\end{eqnarray*}

Substituting into $\omega$ we have:
\begin{eqnarray*}
\lefteqn{2\omega_{(id_G,\theta_o(M))}(\theta_*(u^H),\theta_*(v^F)) }\\
&&= \left( -\overline{\eta}_1,-\overline{\xi}_2-\Ad_{\overline{M}^{-1}}\overline{\xi}_2 \right)
 + \left( \Ad_{\overline{M}}\overline{\eta}_1-\Ad_{\overline{M}^{-1}}\overline{\eta}_1,\Ad_{F_\theta(M)^{-1}}\tfrac{d}{dt}|_{t=0}F_\theta(L_M(r_M(t)))B(M)^{-1} \right)
\end{eqnarray*}

Since $\eta_1\in\fg_M$, the first term and the third term of the last equation cancel out. Thus we obtain
$$2\omega_{(id_G,\theta_o(M))}(\theta_*(u^H),\theta_*(v^F))=(\overline{\eta}_1,\overline{\xi}_2+\Ad_{\overline{M}^{-1}}\overline{\xi}_2)=2(\overline{\eta}_1,\overline{\xi}_2)$$

A similar calculation gives us
$$2\omega_{(id_G,\theta_o(M))}(\theta_*(u^F),\theta_*(v^H))=-2\omega_{(id_G,\theta_o(M))}(\theta_*(v^H),\theta_*(u^F))=-2(\overline{\xi}_1,\overline{\eta}_2).$$
So $\omega_{(id_G,\theta_o(M))}(\theta_*(u),\theta_*(v))=(\overline{\eta}_1,\overline{\xi}_2)-(\overline{\xi}_1,\overline{\eta}_2)\in\bC$,
comparing to $\omega_{(id_G,M)}(u,v)=(\xi_1,\eta_2)-(\eta_1,\xi_2) \in\bC$, we see that $$\omega_{(id_G,\theta_o(M)}(\theta_*(u),\theta_*(v))=-\overline{\omega_{(id_G,M)}(u,v)}. $$
We conclude that $\theta^*\omega=-\overline{\omega}$ along $\epsilon(\Sigma)$ in $\cZ$.
\end{proof}

\end{document}
